%===================================================================================================
% PREAMBLE
%===================================================================================================
\documentclass[conference, a4paper]{IEEEtran}

\usepackage{lipsum}

\usepackage{cite}
\usepackage{amsmath,amssymb,amsfonts}
\usepackage{algorithmic}
\usepackage{graphicx}
\usepackage{textcomp}
\usepackage{xcolor}

\def\BibTeX{{\rm B\kern-.05em{\sc i\kern-.025em b}\kern-.08em
    T\kern-.1667em\lower.7ex\hbox{E}\kern-.125emX}}

%===================================================================================================
% CONTENT
%===================================================================================================
\begin{document}

\title{Image Deblur and Denoise}

\author{
\IEEEauthorblockN{Cosmin-Eusebiu Patrașcu}
\and
\IEEEauthorblockN{Vlad-Ștefan Moreanu}
}

\maketitle

\begin{abstract}
    Hello world! \cite{beyond_a_gaussian}
\end{abstract}

%===================================================================================================
% INTRODUCTION
%===================================================================================================
\section{State-Of-The-Art}

The landscape of image restoration and synthesis has advanced rapidly,
moving from highly optimized Convolutional Neural Networks (CNNs) to highly efficient Vision Transformers and diffusion models.
Each state-of-the-art (SOTA) model introduces distinct architectural innovations to overcome the trade-offs between computational cost and restoration accuracy

\subsection{Convolutional Neural Networks}

DnCNN (Denoising Convolutional Neural Network) introduces a residual learning strategy combined with batch normalization.
Instead of predicting the clean image directly, the model learns to predict the residual noise.
This speeds up training, stabilizes the network and allows a single model (DnCNN-B) to handle blind gaussian denoising. \cite{beyond_a_gaussian}

This strategy allowed DnCNN-B to outperform AR-CNN, a specialized model, by 0.3 dB in PSNR.

FFDNet brings a tunable noise level map as a direct input, allowing a single network to flexibly handle complex, non-uniform noise.
Additionally, it processes the image in a downsampled sub-image space, which vastly improves inference speed and expands the receptive field without losing modeling capacity. \cite{FFDNet}

This model can keep a performance of 26-27 dB PSNR across a vast noise level range, while being significantly faster then DnCNN.

\subsection{High-Efficiency and Adversarial Architectures}

DeblurGAN approaches blind motion deblurring as an image-to-image translation task using Conditional Generative Adversarial Networks (cGANs).
This introduces a multi-component loss function, combining WGAN-GP and perceptual loss, to generate sharp and realistic textures
rather than the overly smooth results found in basic CNNs. It's efficient architercture makes it 5 times faster than other models,
making it a good choice for downstream object detectuin tasks. 
On the GoPro dataset, it achieved a PSNR of 29.7 dB and an SSIM of 0.958. \cite{Deblurgan}

MIMO-UNet brings the re-engineering of the coarse-to-fine strategy into a single U-Net.
It uses a Multi-Input Single Encoder (MISE) to process multi-scale blurry images simultaneously,
and a Multi-Output Single Decoder (MOSD) to output multiple deblurred images,
achieving SOTA accuracy while running over 4 times faster than previous complex models. \cite{Rethinking}

The largest model (MIMO-Unet++) has a PSNR of 32.68 dB on the GoPro dataset, with a processing time of 0.004 seconds.

\subsection{Vision Transformers}

SwinIR, based on the Swin transformer, replaces traditional convolution layers with Residual Swin Transformer Blocks (RSTB).
This has the advantage of local attention processing and shifted-window mechanisms for modeling long range dependencies.
The model. For the super-resolution task it has achieved a PSNR of 32.7 dB, while having significantly less parameters than models with the same results. \cite{Swinir}

Uformer brings transformer capabilities into a classic U-shaped arhitecture.
It's contribution is the Locally-enhaced Window (LeWin) Transformer block,
which calculates self-attention only within non-overlapping windows to reduce computational complexity.
It also introduces a learnable multi-scale restoration modulator to handle varied image degradations. \cite{Uformer}

On image denoising, Uformer achieved 39.9 dB in PSNR, outperforming heavier models.
For debluring, on the GoPro dataset, the model obtained a result of 33 dB 

Restormer brings Multi-Dconv Head Transposed Attention (MDTA), 
which calculates self-attention across the feature (channel) dimension rather than the spatial dimension,
bringing the self-attention complexity down to a linear scale, which is quadratic for normal transformers.
It also introduces a Gated-Dconv Feed-Forward Network (GDFN) to act as a filter for less important features.
For the task of denoising, it has achived a PSNR of 40 dB, and 33 dB for deblurring. \cite{Restormer}

\subsection{Generative Synthesis}

While diffusion models are powerful, calculating them directly in pixel space is incredibly slow and expensive.
Latent Diffusion Models (LDM) bring the innovation of decoupling perceptual and semantic compression.
By first training an autoencoder to compress images into a lower-dimensional latent space, the diffusion process can occur much more efficiently.
This allows the model to strike a near-optimal balance between complexity reduction and detail preservation.
Furthermore, LDMs introduce cross-attention layers into the diffusion UNet, making them incredibly flexible for multi-modal conditioning. \cite{High-Resolution_image}


%===================================================================================================
% 
%===================================================================================================
\section{Method and Dataset Description}


\lipsum[1-3]

%===================================================================================================
% 
%===================================================================================================
\section{Experiments and Results}

\lipsum[1-3]

%===================================================================================================
% CONCLUSIONS
%===================================================================================================
\section{Conclusions}

\lipsum[1]

%===================================================================================================
% BILBIOGRAPHY
%===================================================================================================
\bibliographystyle{IEEEtran}
\bibliography{references}



\end{document}